% Standalone assessment. All references are inline.
% Intentionally uncompiled, as requested. No external inputs are required.
\documentclass[11pt]{article}
\usepackage[T1]{fontenc}
\usepackage[utf8]{inputenc}
\usepackage{lmodern}
\usepackage[margin=1in]{geometry}
\usepackage{amsmath,amssymb}
\usepackage{enumitem}
\usepackage{microtype}
\usepackage[hidelinks]{hyperref}
\setlist{itemsep=3pt,topsep=5pt,leftmargin=*}
\setlength{\parskip}{3pt}
\title{Reevaluating the categorical phenomenology research contract}
\author{}
\date{4 November 2026}

\begin{document}
\maketitle

\begin{abstract}
The contract should continue to demand genuine construction, substantive comparison, explicit warrant, and precise boundaries. Its principal weakness is that these safeguards are better developed than the rules for choosing research that bears on the larger question. A succession of correct, carefully qualified local results can satisfy much of its working discipline while leaving the formation of stable mathematical and scientific theories almost untouched. The current text already protects broad generative and comparative ambitions, allows heterogeneous mechanisms, and makes its native-realization tests conditional. Those protections should be acknowledged and enforced, not rediscovered through a caricature of the contract.

The provisional recommendation is to distinguish a stable programme question from revisable local research commitments, and to assess the explanatory contribution of each local result before letting it determine the next problem. The programme concerns how observations, mathematical evidence, and admissible operations give rise to useful structures and theories; what unifies genuinely different constructions; what transports between contexts; and which differences must remain. Its diagnostic examples cannot define its eventual scope. Recent work has obtained real comparison, obstruction, acquisition, and language-change results. It has not yet established the broader account. This document evaluates the contract and possible changes; it does not replace the contract, adopt new obligations, or initiate a research programme.
\end{abstract}

\tableofcontents

\section{The question being reevaluated}
\label{sec:question}

The immediate concern is not whether the recent manuscripts contain correct mathematics. It is whether the contract keeps the research directed towards its prize when there are many mathematically legitimate next lemmas. The user's formulation is particularly clear: ``A contract and its obligations ought to keep the prize in view.'' The accompanying correction is that the map is not exhaustive and does not explain how stable mathematical or scientific theories arise. Questions of unification, transport, and reuse across contexts belong to that ambition, but reuse is not its complete definition.

This assessment reads the actual amended contract \cite{Contract} against the original method map \cite{Map}, the contextual-realization problem note \cite{ProblemNote}, and the later constructions discussed in Section~\ref{sec:results}. The contract's claims about its earlier research ledger are used as its stated baseline, not as a fresh verification of every earlier manuscript. For the later notes, the assessment examines their theorem statements, arguments, and explicit scope boundaries. It does not certify originality against the entire external literature or independently reprove every result.

Three judgments must therefore remain separate:
\begin{enumerate}
\item whether a result is mathematically valid at its stated scope;
\item whether it explains something important about the programme's question;
\item whether it should determine the next allocation of research effort.
\end{enumerate}
A positive answer to the first does not settle the other two. Conversely, a result can contribute a useful part of the explanation without constituting a new theorem or completing a central commitment. The contract already recognizes several of these distinctions. The issue is whether they are sufficiently effective in directing the work.

\section{What the prize includes}
\label{sec:prize}

\subsection{Formation as well as reconstruction}

The broad question is how structured understanding becomes possible and warranted from observational interaction and other admissible evidence. Depending on the setting, the resulting structure might be a concept system, an algebra of operations, a predictive state, a topology, an explanatory law, a useful representation, a collection of possible models, or a mathematical theory with new objects and relations. These are examples of possible outputs, not an exhaustive list or a sequence every project must traverse.

The distinction between reconstructing an object in a supplied language and forming or revising the language itself is central. A theorem can recover the least sufficient state for fixed experiments while leaving the choice of experiments, significance of the state variables, and formulation of the theory untouched. That theorem is valuable. It does not by itself explain how a stable body of concepts and laws arises.

For scientific theory formation, the missing work can concern why a family of observations supports a particular explanatory structure, which alternatives remain, how new experiments revise the structure, and why its use extends beyond the observations that suggested it. For mathematical theory formation, it can concern the introduction of concepts, sorts, operations, axioms, and proof-relevant distinctions; their interpretation in existing mathematics; and the consequences made available by the new organization. Mathematical evidence here need not be a physical sample. Constructions, examples, counterexamples, proofs, and constraints on representation can all matter. Neither setting requires an assumption-free invention of a theory from an uninterpreted stream.

``Stable'' also needs a local meaning when a theorem is claimed. It might mean invariance under a specified translation, persistence under admissible evidence extension, robustness to a perturbation, or continued explanatory use across a range of contexts. These are different properties. Stability should not be identified with freezing a vocabulary, preserving every old construction literally, or discovering a uniquely true hidden ontology.

There is a serious objection to foregrounding this broader prize: ``theory formation'' can become so elastic that no result counts as enough and no proposed inquiry can be rejected. The remedy is to specify the selected local claim, not to demand a full theory of theory formation from each project. A theorem that meets its stated commitment remains a success at that scope. A later judgment that the programme needs a different direction must not move that theorem's goalposts retrospectively.

\subsection{Unification with differences retained}

A successful common account could be a shared construction, a family of related constructions, a transport principle, a theory of how choices change the output, or a precise explanation of an unavoidable division. The important issue is what the account enables us to derive or understand. A common ambient category and a list of translations are insufficient if they add no such consequence.

There is also no reason to require all constructions to produce the same kind of object. A concept lattice, a minimal automaton, a statistical experiment, and a classifying object can have different universal properties and answer different questions. Unification can explain those differences and still produce substantive shared mathematics. Directed comparisons, partial translations, resource-sensitive transformations, and obstruction results can be as important as equivalences, provided their interpretations and limits are proved.

This is already visible in the original materials. The problem note separates identification from completion, gives consumer-sensitive quotients, requires cross-domain theorem transfer, and explicitly permits a common observational base with doctrine-specific modules if one ambient notion becomes vacuous \cite[sections on success conditions, transport, and the proposed main conjecture]{ProblemNote}. The method map includes scientific model discovery, type formation, categorical logic, inverse problems, active observation, quantitative reconstruction, and many other directions. Its opening explicitly makes the entries provisional and non-exhaustive; its final dependency sketch is expressly not a mandatory pipeline \cite{Map}.

The original materials should themselves remain revisable. Their contextual-realization schema is a serious candidate organization, not a proof that all future theory formation must have that shape. The map's note about institution theory is also explicit: the earlier aversion was largely an artefact of preventing overemphasis. Institutions or other logical tools should be evaluated by what they contribute, not excluded by an accidental historical preference.

Neither transfer nor reuse should replace that schema as a new master criterion. They are important questions alongside formation, warrant, stabilization, and explanatory organization. A result may matter because it changes our understanding of one of these questions before any transport theorem is available. The programme must eventually explain how its achievements relate, but not every achievement has to take the form of reuse across contexts.

\subsection{Tests diagnose the programme rather than constitute it}
\label{sec:diagnostics}

It helps to distinguish three uses of ``test'': an observational probe inside a mathematical problem; a recovery example used to check a proposed construction; and a diagnostic question about the research programme. A theory may make its observational probes part of the formal data without making a fixed recovery suite the definition of phenomenology.

FCA, automata, topology, learning dynamics, and native stochastic realizations can reveal important errors or support a proposed common law. They cannot delimit everything the programme is allowed to become. Passing all presently listed examples would still leave the question of what has been explained. Failing an irrelevant example need not invalidate a more limited theorem. Encountering a genuinely new form of theory formation may require revising the diagnostic suite itself.

The alternative is not a programme immune to failure. Claims still need explicit hypotheses, constructions, comparisons, and counterexamples. The distinction is between testing a claim at its advertised scope and letting the accumulated tests choose the subject of the inquiry.

\section{What the current contract gets right}
\label{sec:protections}

The contract is not a definition of phenomenology as finite stochastic modeling. Its central architecture is significantly broader, and several criticisms would be inaccurate if directed at its literal text.

\paragraph{It protects two positive ambitions.}
The opening and the comparative and generative targets require both an explanation of relations among unlike reconstruction methods and an account of structure not supplied as that output in advance. Comparison must do more than list algorithms; generation must do more than relabel architecture. Either branch can remain incomplete while the other advances. The document explicitly says that additional downstream manuscripts cannot compensate for failure of a central branch \cite[sections on the proposed direction, central targets, and decision gates]{Contract}.

\paragraph{It permits broad and noncanonical forms of emergence.}
Objects, states, operations, relations, geometric or algebraic invariants, dynamic states, hypothesis languages, and families of explanations are all admitted. Useful selection within a supplied search space is allowed if the theorem explains the selection without pretending to have discovered the search space. Exact, nonalgorithmic construction is not excluded. A decomposition is a demanding flagship, not the definition of emergence \cite[the generative target and the flagship subsection]{Contract}.

\paragraph{It does not select a universal categorical language in advance.}
Two-sided presentation is a candidate. One-sided, temporal, approximate, dynamically indexed, and other treatments remain permissible. The relative-computation bridge is explicitly conditional. The contract distinguishes a failed functor from an impossibility of comparison, and non-equivalence from absence of a useful directed map \cite[the proposed direction, comparison semantics, and relative-computation bridge]{Contract}.

\paragraph{It carefully protects epistemic and mathematical meaning.}
Available records, accessible experiments, complete ideal laws, and hypothetical enlarged laws are separated. Forced structure is distinguished from chosen structure. Native realization is distinguished from abstract completion. Equality of predictions is distinguished from intervention compatibility; evidence uncertainty from native nonuniqueness; existence from specified-endpoint and universal lifting; exact identification from stable recovery. These distinctions directly serve the prize by preventing an apparent explanation from assuming what it claims to explain \cite[sections on relative data, acquisition, and visible assumptions]{Contract}.

\paragraph{It already allows local credit and conditional burdens.}
The decision gates say that the entire native and empirical programme need not be completed before an intermediate result is credited. The stress-test introduction says that only controls relevant to the claimed semantics need apply. A pure incidence theorem need not solve feedback. An exact-law theorem need not supply a sampling guarantee. Native Gate D is tied to the interface flagship; another generative theorem can pass its own gate while that flagship remains open \cite[mandatory stress tests and Gates B--E]{Contract}.

These are not ornamental caveats. They are material constraints on an honest reading. Making every theorem solve stochastic assembly, treating every obstruction as a failure of the whole programme, or refusing to credit an exact construction until it becomes a practical learner would be failures of execution, not requirements of the contract.

\section{Where the contract can still lead into the weeds}
\label{sec:diagnosis}

The strongest diagnosis is an imbalance between \emph{validity controls} and \emph{research selection}. The contract is good at identifying overclaims and distinguishing completed from open mathematics. It is less decisive about why one true, new, manageable theorem is a better next step than another. The following mechanisms are supported by the text. They are plausible incentives, not a causal proof that the wording alone produced the recent trajectory.

\subsection{A tractable residual can become the default unit of progress}

Gate A passes ``a tractable theorem-sized question with a genuine residual and explicit controls.'' It also requires a comparison of mechanisms and organizing languages, so it is not permission to choose an arbitrary model. Nevertheless, a precise residual beyond an existing theorem is often easier to specify than its contribution to theory formation. A long sequence of local residuals can remain valid, nontrivial, and carefully audited while never resolving the reason those constructions were being compared.

The dependency diagram places cross-instance tests after comparative or generative positive results. The generality section postpones a materially different later application until the central mechanism is clearer. This is a reasonable way to avoid premature generalization. It also creates a feedback problem: the familiar examples can determine the mechanism, and the mechanism can then determine which unfamiliar examples seem admissible. Heterogeneity enters late enough to become confirmation rather than a source of conceptual change.

The correction is not to require every local proof to cover two or ten domains. It is to make the reason for selecting that proof visible before committing to it: which unresolved phenomenon does it address, which explanatory alternatives could it distinguish, and what would its success change outside the calculation itself? A potential transfer can still be conjectural. A carefully chosen obstruction can be decisive without already carrying a second full instance.

\subsection{The prominent learning target explicitly addresses a smaller question}
\label{sec:learning-gap}

The learning section is unusually candid. After supplying a parameter space, response map, and update family, it distinguishes sufficient-state reconstruction from the formation or selection of a representation during training. It then states that the exact target primarily addresses the former, while the optional empirical branch can address the latter \cite[Learning dynamics as a candidate generative test]{Contract}.

This is an actual limitation of the proposed target, not a strawman. The target can succeed exactly as written by deriving a tractable state sufficient for a supplied learning system. That need not explain how a scientific vocabulary, explanatory law, or mathematical organization comes to be formed. Repeatedly choosing this target can displace the broader question without any false mathematical claim.

The current breadth of the generative target makes theory formation permissible, but permission is weaker than an operative research commitment. If stable theory formation is part of the prize, its place should be explicit rather than represented by sufficient learning state as a proxy. That does not require neural experiments or unrestricted automated discovery. A pure theorem about evidence-relative language change, concept introduction, or the relation between alternative explanatory organizations could address it directly.

\subsection{Visible assumptions do not yet explain their selection}

The supplied/derived ledger prevents a theorem from claiming to generate its own assumptions. It does not explain why a particular hypothesis language, closure doctrine, model class, equivalence, or notion of usefulness was chosen. A collection of impeccable ledgers can therefore produce a collection of separate conditional theories.

A common account needs some positive control over these choices: for example, a comparison explaining what changes when the doctrine changes; a result identifying which conclusions are insensitive to that change; a witness that additional evidence forces revision; or an explanation of why two seemingly similar questions require different output structures. No local result needs to derive all of its own semantics. But the programme cannot leave every important choice external forever and then claim to explain theory formation.

\subsection{The most detailed branch has the strongest pull}

The native targets specify stochastic intermediate links, sequential reuse, interventions, independently assembled routes, finite certificates, and matching resource bounds. The relative-computation and learning branches supply further precise obligations. Such detail creates an actionable backlog. The broader formation question has less comparably concrete guidance, even though it is central to the motivation.

The native demands are largely conditional already, so the primary remedy is to enforce their conditional status and resist treating the longest backlog as the default agenda. A further editorial separation could help: maintain the programme's central questions visibly apart from the exact obligations activated by a particular native claim. Once an interface theorem is claimed, its genuine native obligations cannot simply be discarded. Conversely, an unrelated theory-formation theorem should not inherit them because the document contains them.

One wording issue deserves particular clarification. The cylinder section moves from supported-instance qualifications to ``The new theory must'' supply native and normalization comparisons, whereas the relative-computation subsection is explicitly conditional. The broader qualifications suggest a claim-relative reading, but the scope of this imperative is not fully clear for a non-native, logical, or topological result. A reevaluation should resolve that ambiguity explicitly rather than silently impose or waive the comparison.

\subsection{Novelty and technical strength can crowd out explanatory value}

The strongest-source checks are necessary. Classical constructions and previous achievements must not be claimed again. Yet the repeated requirement for a new residual can make a small new complication seem more valuable than a substantial synthesis of known mathematics that changes our understanding of the programme.

The right distinction is not between novelty and rigor. It is among mathematical correctness, source novelty, explanatory contribution, and programme completion. A synthesis may earn real explanatory credit without earning a priority claim. A new technical theorem may earn mathematical credit without becoming a central advance. The promised new positive consequences remain genuine obligations; crediting the synthesis does not erase them.

\subsection{The stop rules can leave an unproductive trajectory running}

The contract stops claims refuted by counterexamples, reduced entirely to prior work, or made circular by their assumptions. It also forbids turning numerical failure into an impossibility proof. Those rules should remain.

They do not by themselves stop a sequence of correct new refinements whose relevance to the prize is decreasing. A separate prioritization judgment is needed: after a substantial local development, what is now better understood about structure formation, comparison, or transport? If only the local model has become more elaborate, the result should be retained at its proper scope and the next extension reconsidered. This is a reason to redirect attention, not to retract valid mathematics or declare the broader ambition impossible.

\section{What the recent results actually contribute}
\label{sec:results}

The recent work is not uniformly peripheral. Several results supply exactly the kind of distinction or bridge a general account would need. Their contribution should be stated positively and then bounded. Calling all of them ``the weeds'' would obscure both the useful mathematics and the actual reason the broader question remains open.

\subsection{A genuine bridge among unlike native constructions}
\label{sec:common-result}

\emph{One operational construction for residual automata, formal concepts, and topology} does more than run three algorithms on the same incidence table. Compatible sequential operations lift canonically to the concept lattice, and a closure comparison intertwines the supported concept structure with the Dowker construction on their exact partial-action domains. The theorem also proves the domain law for composition. Total state dynamics need not give a total action on simplices \cite[the theorem One coherent partial action in the two native languages]{CommonGeneration}.

This is explanatory leverage: it identifies the retained information that makes a common action possible, locates the correct partiality, and shows why stronger identifications fail. It supports a form of unification that preserves different output types rather than erasing them. It therefore answers part of ``what unifies?'' and ``what transports?''

Its main finite setting is deterministic and Boolean, with a completed operational presentation; the note also states a qualified set-sized extension. It does not discover the observation interface, reconstruct an unknown scientific law from samples, or make the three unmarked outputs equivalent. Its attribution section presents the construction as a proved synthesis without a priority claim. Those limits do not cancel the bridge. They do block its promotion to a general theory of formation.

\subsection{Preservation results expose a choice of meaning}
\label{sec:generated-result}

\emph{Generated Observations and Conservative Characterization Formation} gives exact simultaneous criteria for preserving FCA closure, future-response equivalence, and the literal Dowker complex under test extension. It constructs the largest extensional completion preserving those structures. It also shows why definability, observational availability, and admission as a topological witness are different matters \cite[Three extension criteria, Largest simultaneous completion, and Theory extension beyond a fixed incidence domain]{Generated}.

For example, on two objects with singleton test extents, admitting their disjunction as a new geometric witness fills the edge. The old two-point complex becomes contractible even though the logical expression is entirely defined from existing predicates. The lesson is not that disjunction should be prohibited. It is that preservation of a chosen incidence topology is not an automatic consequence of logical conservativity, and that the role of a test must be specified.

This makes the choice of characterization mathematically consequential. It is useful foundational evidence for the programme. However, a largest conservative completion is deliberately designed to preserve selected existing structures. It does not select the scientifically useful structures, justify preserving them, or explain when a theory should change them. Turning conservativity into the overarching objective would lose the generative ambition. The note itself explicitly identifies its fixed-domain boundary and discusses signature change, model morphisms, proof structure, and classifying constructions as additional questions.

\subsection{Acquisition can target structure rather than every unknown datum}
\label{sec:evidence-result}

The finite evidence theorem distinguishes uncertainty about a relation from uncertainty about its reconstruction. Structural conflicts can be detected through columns and rows without enumerating every concept or simplex. Licensed experiments must separate worlds with different structural targets; they need not identify every aspect of the world. In the four-world example, one experiment settles the FCA and Dowker targets while leaving part of the labeled table undetermined. Two legitimate FCA uncertainty constructions also vary in opposite directions as evidence improves \cite[Exact finite target identification; Structural certificates and the cost of useful queries; Two-sided closure reconstruction and refinement; the four-world example]{Evidence}.

The executable-experiment extension adds another useful distinction. In its four-state parity family, only two structural alternatives remain, but exact worst-case identification requires $r$ terminal binary readings when there are $r$ independently parameterized controls. A positive answer can be a short certificate without being easy to find. That exact lower bound is not a bounded-error statistical lower bound, and the control alphabet grows with $r$ \cite[the four-state family and Accessible reconstruction versus an inaccessible target]{Evidence}.

The sequential-acquisition result then explains why separate availability of experiments is insufficient when executing one changes what can be done next. A belief-state reachability criterion supplies finite positive and negative certificates for identifying the specified invariant under deterministic finite experimental dynamics. More sharply, two individually recoverable characterizations need not have a recoverable common refinement under destructive access: there need not be a single finest operationally recoverable quotient. A world-faithful replay mechanism restores the static separation criterion for a finite model family and a globally admissible replayable menu \cite[Sequential invariant-identification criterion; Failure of common-refinement closure; Replay restores the static separation criterion]{Evidence}.

This is a conceptual qualification of the original largest-identification ambition. It does not refute the contextual quotient theorem for a supplied semantic test family. It shows that the semantic quotient and the information jointly obtainable by an executable policy are different objects. Their agreement requires access hypotheses. The consequence is relevant beyond the chosen FCA or Dowker target: what counts as adequate evidence depends jointly on the target and the observation process.

These results still assume a candidate family, interpreted experiments, and a target. They do not validate those assumptions or generate the governing interpretation. Their exact finite criteria draw on established identification and control ideas. Their value here is the explicit structural consequence and separation of previously conflated questions, not an asserted new universal learning principle.

\subsection{A positive source law is identified rather than inserted}
\label{sec:two-stage-result}

The two-stage characterization note deserves separate credit for actual generation. Its first stage supplies a task-preserving quotient and compatible response geometry. Under independently stated normalization and radius assumptions, and with admitted simultaneous four-patch access, saturation of the fourth-moment bounds then forces an initially unknown source law to have origin-plus-signed-axes support with determined weights. The result constructs a positive output generator and predicts sixth-order consequences not used in identification \cite[Radius saturation, The whole-family output generator, and A held out sixth order consequence]{TwoStage}.

This is not merely preservation of a supplied source distribution. It is conditional structural identification with a further predictive consequence, and therefore direct evidence of generative progress. Its scope is narrow: the affine shared-source family, radius bound, saturation regime, positivity, and access assumptions are substantial; a mathematical chart does not create physical access. The held-out consequence is a prediction, not a report of newly collected validation data. The source credits classical extremal-moment and tensor methods and makes no priority claim. It neither identifies a unique microscopic ontology nor explains the formation of its surrounding scientific theory.

\subsection{Language revision reaches beyond adding another state}
\label{sec:language-result}

The capacity and characterization-revision notes are especially relevant to the concern about specific models. They distinguish insufficient capacity within a language from inadequacy of the language's operations or readout. The sequence $1/(n+1)$ has infinite Hankel rank, yet it has a one-real-state rational recurrence
\[
 x_{n+1}=\frac{x_n}{1+x_n},\qquad x_0=1.
\]
The reciprocal coordinate makes the transition affine while making the readout reciprocal. This is a coordinate transport inside an appropriate rational language; replacing all finite linear weighted automata by that language is a different, expressive change \cite[Exact and approximate expressive separation]{LanguageRevision}.

The later results connect certified inadequacy and assumption-relative candidate identification to a realized language change, with explicit error comparisons and limits on the replacement dynamics. The capacity note separately shows why an unconstrained rank increase is not automatically an automaton repair, and why reverse-mode compatibility is not optimizer commutation \cite{Capacity,LanguageRevision}.

The explanatory gain is the separation of four possible repairs: more parameters, different coordinates, different operations, and a different semantic task. This directly bears on theory formation. The running example and candidate rational language remain supplied; finite observations do not refute the union of all finite WFA sizes, and the exact separation does not establish a poor approximation bound at every nonzero tolerance. The notes therefore do not yet give a general method for inventing or selecting the right language.

\subsection{A scientific example separates exact information from usable warrant}

The oscillator example shows that complete passive observations on a fixed time grid can leave the sign of an initial velocity unresolved within the stated oscillator family. An ideal exact off-grid observation can separate aliases, while finitely many noisy passive observations need not remove unbounded aliasing. A calibrated intervention can repair the particular identification problem under different explicit assumptions \cite[the scientific example with aliasing and calibrated intervention]{Evidence}.

This explains why more of the same data, exact mathematical distinguishability, and robust experimental access are different. It does not infer the oscillator law itself. Its correct programme role is an example that tests and motivates a distinction about warrant, not a substitute for a theory of scientific law formation.

\subsection{The overall assessment of the recent trajectory}

Taken together, these results establish more than technical bookkeeping. They give a genuine operational bridge, distinguish incompatible meanings of preservation, make structural uncertainty experimentally actionable, identify a conditional positive source law with further predictions, and explain why evidence can require a change of language rather than more capacity. These are credible pieces of a broader theory.

The unclosed loop is their integration into an account of formation and comparison beyond their chosen presentations. The current sequence has largely answered what follows once a domain, observation grammar, reconstruction target, or candidate replacement is specified. It has less often explained why that organization is appropriate, how it becomes stable, or how the resulting law informs a genuinely different setting. The bridge in Section~\ref{sec:common-result} is a real exception to a purely local trajectory, but it is not yet the whole answer.

The appropriate response is therefore to identify and reuse the conceptual yield of these results, not simply discard them or demand more complications inside each example. Nor should the contract's older ``current gaps'' ledger be treated as permanently current. Later results must be credited against specific obligations, while completion and novelty judgments remain separate.

\section{Which gates should be retained and which should change}
\label{sec:gates}

The following are provisional recommendations about the existing gates. They are not replacement clauses. The main changes concern selection, the distribution of burdens across projects, and when synthesis is assessed.

\subsection{Gate A should select for relevance as well as a new residual}

\textbf{Retain:} a theorem-sized question, explicit semantics, an honest supplied/derived boundary, the strongest relevant predecessor, noncircular assumptions, and a clear possible failure.

\textbf{Revise:} before the local carrier or favorite formalism becomes fixed, state the observation-to-structure phenomenon that needs explanation and why this particular result would change the programme's answer. A prospective transfer, a comparison of competing explanations, or an important obstruction can supply that reason. It need not be another completed theorem at entry. The original Gate A prescribes a broad source comparison for the first research problem; it should not be misread as a fresh omnibus audit for every later lemma. Narrowing that original initial comparison for a different first target would be a proposed amendment, not an obligation silently treated as optional.

\textbf{Reason:} the bridges and separations in Section~\ref{sec:results} are valuable because of their consequences, not simply because they leave a new lemma beyond the preceding paper. This change would distinguish legitimate mathematics from the best next programme step without confusing the two.

\subsection{Gate B should distinguish kinds of generation}

\textbf{Retain:} genuinely derived output, independent warrant, visible supplied structure, counterfactual dependence, ambiguity, and a positive new construction on a nontrivial family. Original Gate B itself makes that positive commitment. An obstruction-only local result can earn credit without passing it.

\textbf{Clarify:} distinguish construction within a supplied language, selection of a structure within a supplied search space, and formation or revision of the language and its interpretation. None should silently stand for another. The learning target in Section~\ref{sec:learning-gap} belongs primarily to the first kind unless additional formation results are proved.

\textbf{Reconsider:} the requirement that every generative success already include a substantive relation to another reconstruction mechanism. That relation is appropriate when the local theorem advertises unification; otherwise it can remain an explicit programme obligation. A genuinely important formation theorem should not be distorted to manufacture an immediate cross-domain bridge. This is a real possible change to the acceptance architecture, not merely enforcement of existing wording.

The reciprocal danger is endless local generation with deferred comparison. Any relaxation therefore needs the synthesis judgment in Gate F below. It must not retroactively mark the paired programme complete.

Except for the explicitly reconsidered local comparison requirement, the six original nonvacuity requirements and the new-result and ambiguity-boundary obligations are retained by this recommendation. Their abbreviated discussion here does not weaken them. This assessment is not a complete pass/fail audit of every existing result against Gate B.

Two safeguards should remain claim-relative. An independent hypothesis need not always be a finite algorithmic test: compactness, density, regularity, or a stated exact closure property can be legitimate premises of an abstract theorem, while effectiveness is an additional claim. Likewise, retaining genuine synthesis difficulty is a useful control against making a construction problem trivial; difficulty is not a universal measure of conceptual value. A simple consequence can reveal an important relation, and a difficult calculation can remain peripheral. This distinction allows contribution credit; it does not waive the current Gate B's nontrivial positive-family requirement.

\subsection{Gate C should preserve strong comparison without fixing the landscape}

\textbf{Retain:} independent interpretations of the compared structures; the promised new positive bridge; exact direction, loss, and obstruction semantics; and a consequence beyond a common wrapper or preservation by uniqueness alone. Original Gate C has a positive commitment, so a non-comparison theorem alone does not pass it.

\textbf{Clarify the scope and consider revision separately:} the incidence, predictive, and native-factorization trio is a valuable current comparison problem, not the constitutive boundary of a complete theory. A later result can earn comparative significance by explaining a materially different contrast. But the originally specified trio remains an obligation of the original central target unless a deliberate amendment replaces it; broadening the field does not discharge that target. Conversely, three familiar cases sharing notation do not pass merely because the prescribed names appear.

The actual partial-action bridge supplies a useful model of the required evidence: it identifies a common operation law and the exact domains on which it survives, while refusing to identify the unmarked outputs. Its programme credit does not settle a priority claim or every native-factorization obligation.

\subsection{Gate D should remain conditional and demanding}

\textbf{Retain:} the native obligations for a claim that an emerged interface supports native use. Stochastic links, sequential reuse, independent assembly, intervention semantics, and resources cannot be omitted after they have been promised.

\textbf{Change the organization rather than weaken the mathematics:} visibly attach each obligation to the claim that activates it. The contract already allows other forms of generation to progress without completing the flagship. Applying that rule consistently may suffice; moving the detailed branch out of the central decision path would make the rule harder to overlook.

A decision to retire or reduce the flagship would be a separate substantive change. This assessment does not take that decision, and a classification or obstruction does not automatically discharge its positive commitments.

\subsection{Gate E should stay strict about evidence but enter the design earlier}

\textbf{Retain:} the separation of full laws, exact queries, samples, hypothetical probes, and executable access; explicit uncertainty and resource assumptions; and the prohibition on calling fitted or selected structure identified.

\textbf{Adjust the sequencing:} when the intended contribution concerns actual formation from evidence, the evidence and access model should shape the question from the beginning, not appear only as a final validation of an exact construction. The sequential and oscillator examples show that the access model can change the problem itself. Exact-law mathematics remains legitimate without an empirical implementation when that is its advertised claim.

This is largely a stronger application of distinctions already present in the contract, rather than a demand that every theorem become an empirical study.

\subsection{Gate F should judge the explanation throughout the programme}

\textbf{Retain:} the requirement that common mathematics yield a consequence beyond separately described baselines, with failed transfers and scope restrictions explicit.

\textbf{Revise:} do not reserve this assessment for a final stage. After a substantial local result, ask what it changes in the account of formation, comparison, or transport before choosing its successor. Some results will be valuable supporting mathematics; some will settle a central relation; some will suggest that the organizing language is wrong. The next question should reflect that judgment.

Evidence of programme success must eventually include substantive generality. It need not mean that every local theorem has the same number of instances, that every domain has been sampled, or that one theorem settles all the original ambitions. The relevant unit is a coherent body of results whose common claims are actually supported.

This does not replace the original requirement to apply the same substantive theorem or law across unlike instances. Its named algebraic and stochastic or positive comparisons, together with a linear treatment or a precise account of the additional analytic assumptions, remain part of the existing generality obligation unless explicitly revised. Coherence among papers alone is insufficient to discharge that obligation.

\section{Alternative ways to organize the commitments}
\label{sec:architectures}

\subsection{Keep the contract and enforce its existing hierarchy}

The least disruptive option is to preserve the current text and enforce its stated distinctions: central positive targets, conditional branches, local credit, unlike-instance transfer, and no manuscript-count substitute. Add a brief programme assessment when selecting the next problem.

This option has a strong defense. Much of the reported drift would already violate the intended hierarchy if a finite stochastic branch were treated as the entire programme. It avoids multiplying documents and moving the goalposts after results arrive. Its weakness is that the detailed backlog and sequencing remain, while formation still lacks an explicit target commensurate with the user's ambition. If the same narrowing recurs under careful execution, enforcement alone is insufficient.

\subsection{A programme question with replaceable local research commitments}
\label{sec:two-level}

This is the strongest provisional option. Keep the broad questions, epistemic distinctions, positive ambitions, and standards of evidence at programme level. For a selected investigation, state a limited question, assumptions, deliverables, and stopping conditions. Assess local success at that scope, then separately assess what it contributes to the programme.

The advantage is that a local theorem need neither carry the entire ambition nor become its substitute. A native interface project can retain its demanding obligations; a mathematical language-formation project can have different ones. Heterogeneous methods remain possible without being forced into an early universal mechanism.

The principal risk is fragmentation. An unlimited series of local contracts can become the same collection of papers with clearer bookkeeping. The safeguard is substantive synthesis: local success never by itself validates the prize, and choosing another project requires an account of the unresolved common question it will address. The programme must remain capable of revising its organizing ideas, not merely its examples.

This should be a brief intellectual judgment at major choices of direction, not a new administrative layer. It does not require a gate for every lemma, a compulsory synthesis manuscript after every result, a fixed cross-domain count, or repeated audits of irrelevant literature. If the reorganization adds paperwork without changing the choice or interpretation of research, it has not addressed the problem.

\subsection{Develop one common mechanism first}

A mechanism-first strategy could concentrate on a contextual realization, closure, factorization, or other common construction, and use examples to determine its hypotheses. It offers cumulative mathematics and can reveal a reusable proof rather than a family resemblance. The original problem note reasonably proposed such a direction.

Its risk is selecting only problems the mechanism can express. A permissive wrapper can appear general while hiding the real work in the instance-specific assumptions. Alternatively, a strong mechanism can wrongly exclude formation that requires new syntax, new interpretation, or different notions of warrant. This option is justified only while the mechanism earns its centrality through consequences and remains revisable under difficult contrasts.

\subsection{Organize by explanatory problems and compare methods afterwards}

A problem-first strategy would begin with questions such as what evidence warrants introducing a new distinction, what survives a change of theoretical language, or when an explanatory organization remains useful across contexts. Different constructions could then compete or contribute to the answer.

This reduces dependence on the current finite models and keeps formation visible. Its risk is a landscape of insightful case studies without common mathematics. The answer is not to force all cases into one formalism, but to require explicit results explaining where constructions agree, where they differ, and what can be transported. A problem-first search can therefore complement the two-level organization rather than replace its rigor.

\subsection{Make classification the main outcome}

A proved classification of irreducibly different mechanisms, with directed translations and exact obstructions, could be a major result. It would avoid forcing a unity the mathematics does not support. But adopting classification as sufficient programme success would change the present independent positive generative commitment. Failure of several candidate bridges is not yet such a classification, and non-equivalent outputs do not establish irreducibly different mechanisms. This option remains legitimate for discussion, not something an audit can adopt by relabeling unmet targets.

\paragraph{Assessment.}
There is no evidence here that one architecture must govern the entire programme permanently. The two-level option, combined with problem-first selection and selective mechanism development, best addresses the present complaint. Merely widening the recovery checklist would not: it could replace a narrow sequence of examples with a larger sequence of examples while leaving the same unanswered question.

\section{Concrete provisional changes in research direction}
\label{sec:recommendations}

The following recommendations concern what a reevaluated contract should make possible and require. They do not select or commission the next theorem.

In compact form: retain the paired aims and core standards; foreground formation and separate conditional branch obligations; assess the explanation at major research pivots. Decide separately whether to relax Gate B's local comparison requirement or replace any part of the original comparative target. Neither change follows automatically from the recommendation to broaden the programme's attention.

\begin{enumerate}
\item \textbf{Keep the paired ambition, but put theory formation explicitly inside it.} State how the programme intends to address the formation and revision of mathematical or scientific concepts and laws, rather than treating sufficient state reconstruction as an adequate proxy. Do not reduce this to automatic law discovery or insist on a unique inferred ontology.

\item \textbf{Recover the explanatory content of the existing results before extending their models.} The useful content includes partial operational transport, distinct roles for generated tests, structure-targeted evidence, executable sequential access, and language change rather than mere capacity growth. State which of these belong to a common account and which remain separate. Their exact hypotheses and different meanings must survive that synthesis.

\item \textbf{Choose the next question by the uncertainty it resolves about the programme.} A promising local continuation should explain a consequential distinction or decide between plausible organizing ideas. The nearest unproved lemma is not automatically the right choice. A genuinely different setting should enter when it can challenge the proposed principle, not merely decorate it after the principle has been fixed.

\item \textbf{Separate proof obligations from diagnostic examples.} Keep rigorous gates for claims actually made. Retain recovery examples and counterexamples as revisable instruments for discovering bad assumptions, hidden structure, and false transfers. Neither the current test suite nor the original map is a completeness specification for phenomenology.

\item \textbf{Keep contribution and novelty judgments separate.} Credit a synthesis when it explains a relation previously left opaque in this programme. Credit a new theorem when its residual survives a focused source audit. Do not confuse either with completion of the broader positive commitments.

\item \textbf{Allow a branch to be correct and still no longer be the priority.} A sequence of increasingly specialized repairs should not perpetuate itself merely because each repair is valid. Retain its mathematical results and open obligations, then reconsider the next step against the prize. Any actual withdrawal of a promised target should be explicit rather than disguised as completion.
\end{enumerate}

Two kinds of investigation would provide useful discrimination without functioning as a new checklist. One would study a change in explanatory language or mathematical presentation where the warranted output includes new objects, relations, or operations, and where interpretation and proof or evidence genuinely constrain the change. The language-revision work is a foothold, but the next contribution would have to concern more than the already supplied rational example. Another would test whether an explanatory law extracted from the recent work yields a nontrivial construction or obstruction in a setting whose semantics differ materially from the finite incidence and automaton presentations. In either case, failure to fit the current organizing language would be informative.

These are directions for choosing a question, not fully specified targets and not a claim that either is already the best next project. The map can supply possible tools and comparison settings; it cannot decide the inquiry in advance.

\section{Judgments that remain open}
\label{sec:open}

Some choices cannot be settled merely by auditing the wording. They should remain visible rather than being silently resolved through the next convenient model.

\paragraph{What counts as the central formation achievement?}
There are meaningful differences among deriving structure from a complete law, learning within a known hypothesis class, revising a characterization language, and explaining the formation of a stable theory. Which of these is the next central objective remains a research-design judgment. The user's broader aim prevents the first from silently exhausting the fourth.

\paragraph{How unified should the outcome be?}
A single mechanism may be possible in part of the landscape. A structured family with strong bridges may be the more faithful result. The evidence presently supports local common laws and explicit boundaries; it does not establish that one universal construction is either possible or impossible. This question should be tested, not settled by preference for one categorical technology.

\paragraph{How much comparison belongs inside a local generative commitment?}
The current contract couples them strongly. Relaxing that local coupling could admit more ambitious formation work, but it also risks indefinite postponement of unification. The two-level option needs an explicit programme-level return to comparison, without imposing a fixed number of domains on every theorem.

\paragraph{Which native ambitions should remain active?}
The unsupplied-interface and multistage targets remain demanding and potentially valuable. Their technical detail is not itself a reason to discard them, and this assessment does not do so. Their priority relative to theory formation and other comparisons should be decided on explanatory grounds rather than inherited from the size of their backlog.

\paragraph{What makes a retained distinction useful?}
Every exact theorem is relative to a target, consumer, or preserved law. The programme should investigate some of the dependence on those choices. It should not promise to derive all relevance from observations alone, nor let relevance remain permanently unexplained while presenting the resulting theory as complete.

\paragraph{Has the contract itself failed, or has its execution failed?}
There is evidence for both a textual asymmetry and a prioritization problem. The current text already forbids many of the narrow interpretations that would explain the drift. Yet its learning target, detailed native branches, and staged generality make some narrower trajectories entirely legitimate. A modest reorganization plus stronger selection may suffice. A wholesale replacement is not yet justified by this assessment.

\section{Conclusion}
\label{sec:conclusion}

The contract should retain its mathematical seriousness. The supplied/derived boundary, source discipline, contextual validity, native semantics where claimed, and distinction between evidence and ideal laws are part of keeping the prize in view. Removing them would make broad claims easier without making the programme more explanatory.

The needed correction is to what selects and connects the work. The recent results show that common constructions and sharp distinctions are possible. They also show how easily a sequence can remain inside supplied languages, finite model classes, and increasingly exact preservation questions. Those results deserve credit at their scope; they do not yet explain the formation of stable mathematical and scientific theories.

The strongest provisional recommendation is therefore to retain the paired comparative and generative ambition, make theory formation explicit, keep native and empirical burdens attached to the claims that activate them, and separate local theorem success from the programme's synthesis judgment. The next problem should be chosen for the explanation it can change, with a concrete model used to investigate that question rather than to define it. Nothing in this assessment adopts those changes or replaces the current contract.

\appendix
\section{Sources and the scope of this assessment}
\label{app:sources}

References below identify the programme documents used in this assessment. The bibliography is inline so that this source is standalone. Titles and theorem names in the text locate the relevant passages without requiring external LaTeX cross-references. The original method map and problem note were read in full, as was the amended contract. Selected later theorem statements, proofs, examples, and scope discussions were inspected for the claims in Section~\ref{sec:results}. Bibliographic novelty claims made or withheld in those notes retain their original status.

This file was prepared as uncompiled \LaTeX{} source. Static source checks cover environments, labels, references, and bibliography keys. No compilation, page-layout inspection, new mathematical experiment, or new external literature review is represented as completed by this assessment.

\begin{thebibliography}{99}

\bibitem{Contract}
\emph{A research contract for categorical phenomenology.}
Revised discussion draft with parametric amendment, dated 3 November 2026.
Source filename: \texttt{Categorical\_Phenomenology\_Parametric\_Amendment.tex}.
The assessment concerns its substantive targets, qualifications, and decision gates; it does not infer or change its acceptance status from the draft's own status notices.

\bibitem{Map}
\emph{Methods for generative categorical phenomenology.}
Original programme method map, including its opening qualification, forty-two thematic sections, and candidate entry points. No authorship date is inferred from retrieval or file metadata.
Recovered attachment: \texttt{categorical-phenomenology-method-map(3).md}.

\bibitem{ProblemNote}
\emph{Categorical Phenomenology: A Problem Note on Contextual Realization.}
Original draft research note. Source reference in the method map:
\texttt{contextual-realization-problem-note.tex}.
Recovered attachment: \texttt{categorical\_phenomenology\_problem\_note.tex}.

\bibitem{CommonGeneration}
\emph{One operational construction for residual automata, formal concepts, and topology.}
Programme note, 4 November 2026.
\texttt{common\_generation\_note.tex}.
Relevant theorem: \emph{One coherent partial action in the two native languages}; relevant discussion: \emph{Scope, limitations, and the contribution of the synthesis}.

\bibitem{Generated}
\emph{Generated Observations and Conservative Characterization Formation.}
Programme note.
\texttt{generated\_characterizations\_note.tex}.
Relevant results: \emph{Three extension criteria}, \emph{Largest simultaneous completion}, and the sections on theory extension and contribution boundaries.

\bibitem{Evidence}
\emph{Evidence Warranted Reconstruction and Experimental Access.}
Integrated programme note, 4 November 2026.
\texttt{evidence\_warranted\_reconstruction.tex}.
Relevant results: structural comparison and conflict witnesses, two-sided closure reconstruction, the four-state parity family, sequential invariant identification, failure of common-refinement closure, replay, and the oscillator example.
The constituent theorem and experiment modules were also inspected.

\bibitem{TwoStage}
\emph{Two Stage Generative Characterization: From conditional quotient discovery to a positive source law and a held out sixth order prediction.}
Programme note, 4 November 2026.
\texttt{Two\_Stage\_Generative\_Characterization.tex}.
Relevant results: radius saturation, positive output generation, and the sixth-order consequence.

\bibitem{Capacity}
\emph{Evidence-Driven Capacity Change: Certificates, Reconstruction, and Live Gradient Extensions.}
Programme note, 4 November 2026.
\texttt{capacity\_expansion\_note.tex}.

\bibitem{LanguageRevision}
\emph{Evidence Driven Revision of Characterization Languages.}
Programme note, 4 November 2026.
\texttt{characterization\_revision\_note.tex}.
Relevant sections: exact and approximate expressive separation, finite evidence, assumption-relative rational identification, realized language change, and the least observable extension.

\end{thebibliography}
\end{document}
